\documentclass[../../main.tex]{subfiles}
\begin{document}
\section{带边流形上的流与流出}\label{section:ch09-带边流形上的流与流出}

一般而言，带边流形上的光滑向量场不一定生成流，因为例如从某些边界点出发的积分曲线可能仅在半开区间上有定义。但流出定理有一个适用于带边流形的变体，它有许多重要应用。

假设 $M$ 是带非空边界的光滑流形。下一个定理描述了从 $\partial M$ 出发的某种"单侧流出"，由在 $\partial M$ 上处处指向内部的向量场所决定。

\begin{theorem}\label{theorem:ch09-24}[边界流出定理]
设 $M$ 是带非空边界的光滑流形，$N$ 是 $M$ 上在 $\partial M$ 每点处指向内部的光滑向量场。存在光滑函数 $\delta\colon\partial M\to\mathbb{R}^+$ 和光滑嵌入 $\Phi\colon P_\delta\to M$，其中 $P_\delta=\{(t;p):p\in\partial M,\;0\leqslant t<\delta(p)\}\subseteq\mathbb{R}\times\partial M$，使得 $\Phi(P_\delta)$ 是 $\partial M$ 的邻域，且对每个 $p\in\partial M$，映射 $t\mapsto\Phi(t;p)$ 是从 $p$ 出发的 $N$ 的积分曲线。
\end{theorem}

\begin{proof}
习题 9-11。
\end{proof}

设 $M$ 是带边光滑流形。$\partial M$ 的一个邻域称为\textbf{领圈邻域}，若它是光滑嵌入 $[0;1)\times\partial M\to M$ 的像，且该嵌入限制为显然的等同 $\{0\}\times\partial M\to\partial M$。

\begin{figure}[ht]
\centering
\includegraphics[width=0.8\textwidth]{fig-09-9.png}
\caption{边界的领邻域}\label{figure:ch09-collar-neighborhood}
\end{figure}

\begin{theorem}\label{theorem:ch09-25}[领圈邻域定理]
若 $M$ 是带非空边界的光滑流形，则 $\partial M$ 有领圈邻域。
\end{theorem}

\begin{proof}
由习题 8-4 的结果，存在光滑向量场 $N\in\mathfrak{X}(M)$，其在 $\partial M$ 上的限制处处指向内部。令 $\delta\colon M\to\mathbb{R}^+$ 和 $\Phi\colon P_\delta\to M$ 如\cref{theorem:ch09-24}，定义映射 $\beta\colon[0;1)\times\partial M\to P_\delta$ 为 $\beta(t;p)=(t\delta(p);p)$。则 $\beta$ 是微分同胚，且在 $\{0\}\times\partial M$ 上限制为恒等；因此映射 $\Phi\circ\beta\colon[0;1)\times\partial M\to M$ 是光滑嵌入，其像开，且限制为通常的等同 $\{0\}\times\partial M\to\partial M$。$\Phi\circ\beta$ 的像即为 $\partial M$ 的领圈邻域。
\end{proof}

领圈邻域定理的第一个应用（其中包括）表明每个带边光滑流形与其内部同伦等价。

\begin{theorem}\label{theorem:ch09-26}
设 $M$ 是带非空边界的光滑流形，$\iota\colon\operatorname{Int}M\hookrightarrow M$ 表示包含。存在正常光滑嵌入 $\rho\colon M\to\operatorname{Int}M$，使得 $\iota\circ\rho\colon M\to M$ 和 $\rho\circ\iota\colon\operatorname{Int}M\to\operatorname{Int}M$ 都光滑同伦于恒等映射。因此 $\iota$ 是同伦等价。
\end{theorem}

\begin{proof}
\cref{theorem:ch09-25}表明 $\partial M$ 在 $M$ 中有领圈邻域 $C$，它是光滑嵌入 $E\colon[0;1)\times\partial M\to M$ 的像，满足对所有 $x\in\partial M$，$E(0;x)=x$。为简化记号，我们将用此嵌入把 $C$ 等同于 $[0;1)\times\partial M$，将 $C$ 中的点记为有序对 $(s;x)$，其中 $s\in[0;1)$，$x\in\partial M$；于是 $(s;x)\in\partial M$ 当且仅当 $s=0$。对任意 $a\in(0;1)$，令 $C(a)=\{(s;x)\in C:0\leqslant t<a\}$，$M(a)=M\smallsetminus C(a)$，后者是 $\operatorname{Int}M$ 中的正则区域。

令 $\eta\colon[0;1)\to[\frac13;1)$ 是递增微分同胚，满足 $\frac23\leqslant s<1$ 时 $\eta(s)=s$，定义 $\rho\colon M\to\operatorname{Int}M$ 为
\begin{align*}
\rho(p)=
\begin{cases}
p, & p\in\operatorname{Int}M\smallsetminus C(\frac23),\\
(\eta(s);x), & p=(s;x)\in C.
\end{cases}
\end{align*}
这两个定义在重叠集 $C\smallsetminus C(\frac23)$ 上都给出恒等映射，故由粘接引理 $\rho$ 光滑。它是到闭子集 $M(\frac13)$ 的微分同胚，故是从 $M$ 到 $\operatorname{Int}M$ 的正常光滑嵌入。

定义 $H\colon M\times I\to M$ 为
\begin{align*}
H(p;t)=
\begin{cases}
p, & p\in\operatorname{Int}M\smallsetminus C(\frac23),\\
(ts+(1-t)\eta(s);x), & p=(s;x)\in C.
\end{cases}
\end{align*}
与前面一样，$H$ 光滑，直接验证表明它是从 $\iota\circ\rho$ 到 $\operatorname{Id}_M$ 的同伦。若 $p\in\operatorname{Int}M$，则对所有 $t\in I$，$H(p;t)\in\operatorname{Int}M$，故 $H$ 在 $(\operatorname{Int}M)\times I$ 上的限制是从 $\rho\circ\iota$ 到 $\operatorname{Id}_{\operatorname{Int}M}$ 的光滑同伦。
\end{proof}

\cref{theorem:ch09-26}是如下 Whitney 逼近定理推广的主要成分。

\begin{theorem}\label{theorem:ch09-27}[带边流形的 Whitney 逼近]
若 $M$ 和 $N$ 是带边光滑流形，则从 $M$ 到 $N$ 的每个连续映射都同伦于一个光滑映射。
\end{theorem}

\begin{proof}
\cref{theorem:ch06-26} 处理 $\partial N=\emptyset$ 的情形，故可设 $\partial N\neq\emptyset$。设 $F\colon M\to N$ 是连续映射，$\iota\colon\operatorname{Int}N\hookrightarrow N$ 是包含，$\rho\colon N\to\operatorname{Int}N$ 是\cref{theorem:ch09-26}中构造的映射，故 $\iota\circ\rho\colon N\to N$ 光滑同伦于 $\operatorname{Id}_N$。\cref{theorem:ch06-26} 表明 $\rho\circ F\colon M\to\operatorname{Int}N$ 同伦于光滑映射 $G$。于是 $\iota\circ G\simeq\iota\circ\rho\circ F\simeq F$，故 $\iota\circ G\colon M\to N$ 是同伦于 $F$ 的光滑映射。
\end{proof}

下一个定理将\cref{theorem:ch06-29}的主要结果推广到映入带边流形的映射的情形。

\begin{theorem}\label{theorem:ch09-28}
假设 $M$ 和 $N$ 是带边或不带边的光滑流形。若 $F,G\colon M\to N$ 是同伦的光滑映射，则它们光滑同伦。
\end{theorem}

\begin{proof}
\cref{theorem:ch06-29} 处理 $\partial N=\emptyset$ 的情形，故可设 $N$ 有非空边界。令 $\iota\colon\operatorname{Int}N\hookrightarrow N$ 和 $\rho\colon N\to\operatorname{Int}N$ 如\cref{theorem:ch09-26}。则 $\rho\circ F$ 和 $\rho\circ G$ 是从 $M$ 到 $\operatorname{Int}N$ 的同伦光滑映射，故\cref{theorem:ch06-29} 表明它们彼此光滑同伦。于是有光滑同伦 $F\simeq\iota\circ\rho\circ F\simeq\iota\circ\rho\circ G\simeq G$。由光滑同伦的传递性（\cref{lemma:ch06-28}），$F$ 光滑同伦于 $G$。
\end{proof}

下一个定理可能是领圈邻域定理最重要的应用。

\begin{theorem}\label{theorem:ch09-29}[沿边界粘合光滑流形]
设 $M$ 和 $N$ 是带非空边界的光滑 $n$-流形，假设 $h\colon\partial N\to\partial M$ 是微分同胚。令 $M\cup_h N$ 为将每个 $x\in\partial N$ 与 $h(x)\in\partial M$ 等同而得到的粘接空间。则 $M\cup_h N$ 是拓扑流形（无边），且具有光滑结构，使得存在正则区域 $M',N'\subseteq M\cup_h N$，分别微分同胚于 $M$ 和 $N$，并满足
\begin{align}
M'\cup N'=M\cup_h N,\qquad M'\cap N'=\partial M'=\partial N'. \label{align:ch09-14}
\end{align}
若 $M$ 和 $N$ 都紧，则 $M\cup_h N$ 紧；若它们都连通，则 $M\cup_h N$ 连通。
\end{theorem}

\begin{figure}[ht]
\centering
\includegraphics[width=0.8\textwidth]{fig-09-10.png}
\caption{沿边界粘合流形}\label{figure:ch09-attaching-manifolds}
\end{figure}

\begin{proof}
为简单起见，令 $X=M\cup_h N$ 表示商空间，$\pi\colon M\sqcup N\to X$ 表示商映射。令 $V\subseteq M$ 和 $W\subseteq N$ 分别为 $\partial M$ 和 $\partial N$ 的领圈邻域，记相应的微分同胚为 $\alpha\colon[0;1)\times\partial M\to V$ 和 $\beta\colon[0;1)\times\partial N\to W$。定义连续映射 $\Phi\colon V\sqcup W\to(-1;1)\times\partial M$ 为
\begin{align*}
\Phi(x)=
\begin{cases}
(-t;p), & x=\alpha(t;p)\in V,\\
(t;h(q)), & x=\beta(t;q)\in W.
\end{cases}
\end{align*}
则 $\Phi$ 在 $V$ 或 $W$ 上的限制是像闭的拓扑嵌入，由此容易推出 $\Phi$ 是闭映射。由于 $\Phi$ 在 $\pi$ 的纤维上为常数，它下降为连续映射 $\widetilde{\Phi}\colon\pi(V\sqcup W)\to(-1;1)\times\partial M$。此映射双射，且因它也是闭映射而是同胚：若 $K\subseteq\pi(V\sqcup W)$ 闭，则 $\pi^{-1}(K)$ 在 $V\sqcup W$ 中闭，因而 $\widetilde{\Phi}(K)=\Phi(\pi^{-1}(K))$ 闭。故 $\pi(V\sqcup W)$ 是拓扑 $n$-流形。另一方面，$\pi$ 在饱和开子集 $\operatorname{Int}M\sqcup\operatorname{Int}N$ 上的限制是单射商映射，因而是到其像的同胚；这表明 $X$ 局部欧氏、维数为 $n$。由于 $X$ 是第二可数开子集 $\pi(\operatorname{Int}M\sqcup\operatorname{Int}N)$ 和 $\pi(V\sqcup W)$ 的并，它第二可数。$M\sqcup N$ 中任意两条纤维可被饱和开子集分离，故 $X$ Hausdorff。因此它是拓扑 $n$-流形。

我们如下定义 $X$ 上的坐标卡集：
\begin{align*}
\bigl(\pi(U),\varphi\circ\pi^{-1}|_{\pi(U)}\bigr),\quad &\text{对 }\operatorname{Int}M\text{ 或 }\operatorname{Int}N\text{ 的每个光滑坐标卡 }(U;\varphi);\\
\bigl(\widetilde{\Phi}^{-1}(U),\varphi\circ\widetilde{\Phi}|_{\widetilde{\Phi}^{-1}(U)}\bigr),\quad &\text{对 }(-1;1)\times\partial M\text{ 的每个光滑坐标卡 }(U;\varphi).
\end{align*}
这些映射是同胚的复合，故在 $X$ 上定义坐标卡；直接验证它们都光滑相容，从而在 $X$ 上定义光滑结构。$\pi$ 在 $M$ 上的限制连续、闭且单射，因而是正常嵌入。在上述任一光滑坐标卡和 $M$ 上相应坐标下，$\pi$ 的坐标表示或是恒等映射或是包含映射，故它是光滑嵌入，其像 $M'$ 因而为 $X$ 中的正则区域。类似考虑适用于 $N$；关系\eqref{align:ch09-14}直接由定义得到。

若 $M$ 和 $N$ 都紧，则 $X$ 是紧集 $M'$ 和 $N'$ 的并，故紧；若它们连通，则 $X$ 是连通集 $M'$ 和 $N'$（有公共点 $\partial M'=\partial N'$）的并，故连通。
\end{proof}

\begin{exercise}\label{exercise:ch09-30}
假设 $M$ 和 $N$ 是带边光滑 $n$-流形，$A\subseteq\partial M$ 和 $B\subseteq\partial N$ 是非空子集，且分别是相应边界的若干连通分支的并，$h\colon B\to A$ 是微分同胚。验证\cref{theorem:ch09-29}的证明只需平凡修改即可通过，表明 $M\cup_h N$ 是带边拓扑流形，并可赋予光滑结构使 $M$ 和 $N$ 分别微分同胚于 $M\cup_h N$ 中的正则区域。
\end{exercise}

\begin{example}\label{example:ch09-31}[连通和]
设 $M_1,M_2$ 是连通的 $n$ 维光滑流形。对 $i=1,2$，令 $U_i$ 是以某点 $p_i\in M_i$ 为中心的正则坐标球，令 $M_i'=M_i\smallsetminus U_i$。习题 5-8 表明每个 $M_i'$ 是带边光滑流形，其边界微分同胚于 $S^{n-1}$。$M_1$ 和 $M_2$ 的光滑连通和，记为 $M_1\# M_2$，是一个光滑流形，通过选取从 $\partial M_1$ 到 $\partial M_2$ 的微分同胚并沿边界粘合 $M_1'$ 和 $M_2'$ 而形成。
\end{example}

\begin{figure}[ht]
\centering
\includegraphics[width=0.8\textwidth]{fig-09-11.png}
\caption{光滑连通和}\label{figure:ch09-smooth-connected-sum}
\end{figure}

若 $M$ 是任意带边光滑流形，\cref{theorem:ch09-26}表明 $M$ 可正常嵌入无边光滑流形（即 $\operatorname{Int}M$ 的一个副本）。下一个例子展示 $M$ 可如此嵌入的另一种方式；此构造的优点是当 $M$ 本身紧时将 $M$ 嵌入紧流形。

\begin{example}\label{example:ch09-32}[带边光滑流形的二重]
设 $M$ 是带边光滑流形。$M$ 的\textbf{二重}是流形 $D(M)=M\cup_{\operatorname{Id}} M$，其中 $\operatorname{Id}\colon\partial M\to\partial M$ 是 $\partial M$ 的恒等映射；它由 $M\sqcup M$ 通过将一份 $M$ 中的每个边界点与另一份中相同的边界点等同而得到。它是无边光滑流形，包含两个微分同胚于 $M$ 的正则区域。容易验证 $D(M)$ 紧当且仅当 $M$ 紧，连通当且仅当 $M$ 连通。（将定义推广到无边流形，当 $\partial M=\emptyset$ 时定义 $D(M)=M\sqcup M$，是有用的。）
\end{example}

虽然带边流形上的向量场不总生成流，但有一种情形可以：当向量场处处与边界相切时。为证明这一点，我们从如下特殊情形开始。

\begin{lemma}\label{lemma:ch09-33}
假设 $M$ 是光滑流形，$D\subseteq M$ 是正则区域。若 $V$ 是 $M$ 上与 $\partial D$ 相切的光滑向量场，则 $V$ 的每条从 $D$ 出发的积分曲线只要有定义就保持在 $D$ 中。
\end{lemma}

\begin{proof}
假设 $\gamma\colon J\to M$ 是 $V$ 的积分曲线，$\gamma(0)\in D$。定义 $T\subseteq J$ 为 $T=\{t\in J:\gamma(t)\in D\}$。我们将证明 $T$ 在 $J$ 中既开又闭；由于 $J$ 是区间，这推出 $T=J$，证明引理。

由于 $D$ 在 $M$ 中闭（由正则区域的定义），由连续性 $T$ 在 $J$ 中闭。为证明它开，假设 $t_0\in T$。若 $\gamma(t_0)\in\operatorname{Int}D$，则由连续性 $t_0$ 的某邻域包含于 $T$，故可设 $\gamma(t_0)\in\partial D$。由于 $V$ 与 $\partial D$ 相切，\cref{proposition:ch08-23}表明存在光滑向量场 $W=V|_{\partial D}$，它与 $V$ 是 $\iota$-相关的，其中 $\iota\colon\partial D\hookrightarrow M$ 是包含。令 $\widetilde{\gamma}\colon(t_0-\varepsilon,t_0+\varepsilon)\to\partial D$ 是 $W$ 的积分曲线，$\widetilde{\gamma}(t_0)=\gamma(t_0)$。由积分曲线的自然性（\cref{proposition:ch09-06}），$\iota\circ\widetilde{\gamma}$ 是具有相同初始条件的 $V$ 的积分曲线，故由唯一性它必等于两者都有定义处的 $\gamma$。这表明在 $t_0$ 的某邻域内 $\gamma(t)\in\partial D\subseteq D$，故如所断言 $T$ 在 $J$ 中开。
\end{proof}

\begin{theorem}\label{theorem:ch09-34}[带边流形上的流]
若 $M$ 是带边光滑流形，$V$ 是 $M$ 上与 $\partial M$ 相切的光滑向量场，则\cref{theorem:ch09-12}的结论仍然成立。
\end{theorem}

\begin{proof}
\cref{example:ch09-32}表明可将 $M$ 视为其二重 $D(M)$ 中的正则区域。由向量场延拓引理，可将 $V$ 延拓为 $D(M)$ 上的光滑向量场 $\widetilde{V}$。令 $\widetilde{\theta}\colon\widetilde{\mathscr{D}}\to D(M)$ 是 $\widetilde{V}$ 的流，令 $\mathscr{D}=\widetilde{\mathscr{D}}\cap(\mathbb{R}\times M)$，$\theta=\widetilde{\theta}|_{\mathscr{D}}$。则\cref{lemma:ch09-33}保证 $\theta$ 将 $\mathscr{D}$ 映入 $M$，其余结论由将\cref{theorem:ch09-12}应用于 $\widetilde{V}$ 得到。
\end{proof}

对带边流形，标准型定理有如下变体。

\begin{theorem}\label{theorem:ch09-35}[边界正则点附近的标准型]
设 $M$ 是带边光滑流形，$V$ 是 $M$ 上与 $\partial M$ 相切的光滑向量场。若 $p\in\partial M$ 是 $V$ 的正则点，则存在 $p$ 某邻域上的光滑边界坐标 $(s^i)$，使得 $V$ 在其中有坐标表示 $\partial/\partial s^1$。
\end{theorem}

\begin{proof}
习题 9-15。
\end{proof}

\end{document}